\documentclass[11pt]{article}

\usepackage{kotex}
\usepackage{geometry}
\usepackage{setspace}
\usepackage{parskip}

\geometry{a4paper,margin=24mm}
\setstretch{1.18}
\setmainhangulfont{AppleMyungjo}
\setsanshangulfont{Apple SD Gothic Neo}

\title{\textbf{프로젝트 기반 자기소개서}\\드림에이스 지원용}
\author{이준영}
\date{2026년 3월}

\begin{document}

\maketitle

\section*{1. 차량 소프트웨어를 기능이 아니라 구조로 본 경험}
현대모비스 부트캠프 PBL에서 저는 구간 정보, 긴급차량 접근, 객체 위험을 하나의 경고 정책으로 통합하는 프로젝트를 수행했습니다. 이 프로젝트는 차량 전체 기준선 위에서 다수 ECU, CAN/Ethernet 네트워크, V2X·ADAS·CGW 논리를 함께 맞춰야 했기 때문에, 기능 하나만 구현하는 방식으로는 정리되지 않았습니다. 저는 그 과정에서 아키텍처와 통신 구조를 정리하고, 핵심 CAPL 로직과 시험 구조를 함께 맞추는 역할을 맡았습니다.

처음에는 기능 정의와 코드만 맞으면 될 것이라고 생각했습니다. 하지만 실제로는 요구사항, 기능 정의, 네트워크 흐름, 시스템 변수, 코드, 시험 기준이 따로 움직이면 구현은 되어도 검증이 닫히지 않았고, 검증이 되어도 그 근거를 설명하기 어려웠습니다. 그래서 무엇을 왜 이렇게 만들었는지, 어떤 경로로 동작하는지, 어떤 기준으로 확인할 것인지를 처음부터 함께 설계해야 한다는 점을 체감했습니다.

이 경험을 통해 저는 차량 소프트웨어를 개별 기능의 집합이 아니라, 구조와 검증 기준까지 포함한 시스템으로 보게 되었습니다. 드림에이스가 IVI, SDV, SILS를 각각 분리된 기능이 아니라 차량 소프트웨어 전체 흐름 안에서 다루고 있다는 점에서, 제가 프로젝트를 통해 쌓은 관점과 맞닿아 있다고 느꼈습니다.

\section*{2. 가상 ECU 기반 개발과 검증을 함께 다뤄 본 경험}
드림에이스 공식 자료를 보며 가장 관심이 갔던 부분은 SILS와 DRIM-SimHub였습니다. 저는 이번 프로젝트에서 CANoe SIL 환경을 사용해 logical ECU 단위로 V2X, ADAS, CGW, TEST 계층을 나누어 동작을 확인했습니다. 실제 차량이 아닌 가상 ECU 기반 환경에서도 기능을 어떻게 설계하고, 어떤 순서로 검증하고, 어떤 결과를 증빙으로 남길지까지 함께 다뤄본 경험이 있습니다.

특히 저는 V2X와 ADAS가 단순한 기능 설명에 머무르지 않고, CAPL 코드, 메시지 흐름, 시스템 변수, 시험 자산으로 이어지도록 정리하는 데 집중했습니다. UT·IT·ST를 분리해 기능 로직 검증과 시스템 기준선 검증을 나누어 보고, 시험 실행 결과는 XML 기반으로 정리해 이후 closeout과 문서 갱신까지 연결했습니다.

이 과정에서 배운 것은 시뮬레이션 환경이 단지 편한 개발 도구가 아니라, 구현과 검증을 더 일찍 연결할 수 있게 해주는 기준선이라는 점이었습니다. 드림에이스가 SILS에서 가상 ECU 기반 개발·검증 플랫폼을 강조하는 이유도 결국 같은 방향이라고 생각합니다. 저는 이미 CANoe SIL 환경에서 비슷한 문제를 다뤄본 경험을 바탕으로, 드림에이스의 SILS와 SDV 검증 업무에도 빠르게 적응할 수 있다고 생각합니다.

\section*{3. 복잡한 시스템을 설명 가능한 결과로 정리한 경험}
이번 프로젝트에서 어려웠던 점은 기능 자체보다, 복잡한 시스템을 팀원과 리뷰어가 같은 기준으로 이해할 수 있게 정리하는 일이었습니다. 구간 경고, V2X, ADAS, 감속 보조, Fail-Safe가 동시에 연결되어 있었기 때문에 작은 표현 차이도 전체 구조를 흔들 수 있었습니다.

저는 이 문제를 코드 수정만으로 풀기보다, 먼저 의미를 정리하는 방식으로 접근했습니다. 어떤 ECU가 어떤 판단을 소유하는지, 어떤 메시지가 어떤 상태로 이어지는지, 시험에서 무엇을 확인해야 하는지를 먼저 정리했고, 그 기준에 맞춰 문서와 구현을 함께 맞췄습니다. 이후에는 시험 실행 결과를 XML 기반으로 정리하고, 엑셀 산출물과 closeout 흐름까지 연결해 반복 작업을 줄이는 방향으로 운영 자동화도 정리했습니다.

이 과정을 통해 시스템 프로젝트에서는 구현 자체만큼, 그 결과를 다시 설명 가능하게 정리하는 일이 중요하다는 점을 배웠습니다. 드림에이스가 공식적으로 `virtualization`과 `openness`를 강조하는 것도 결국 복잡한 차량 소프트웨어를 더 유연하고 설명 가능한 구조로 만들기 위한 방향이라고 이해했습니다. 저는 이런 방향에서 문서, 구조, 검증 결과를 함께 정리하는 역할을 맡고 싶습니다.

\section*{4. 기준을 세우고 팀이 흔들리지 않게 만든 경험}
이번 프로젝트에서 저는 팀장 역할을 맡아 아키텍처, 통신 구조, 전체 방향을 조율하는 일을 했습니다. 프로젝트 범위가 넓다 보니 팀원마다 보는 관점이 달랐고, 구현해야 할 기능도 많았습니다. 이럴 때 제 판단만 앞세우기보다, 팀이 공통으로 따를 수 있는 기준을 먼저 정리하려고 했습니다.

어떤 기능을 먼저 봐야 하는지, 어떤 문서가 기준이 되는지, 어떤 시험이 실제로 중요한지를 계속 정리했고, 회의 때도 ``누가 맞는가''보다 ``무엇을 기준으로 판단할 것인가''를 먼저 맞추려고 했습니다. 그 결과 의견 차이가 있더라도 다시 기준선으로 돌아와 논의할 수 있었고, 프로젝트가 커져도 팀이 쉽게 흔들리지 않았습니다.

드림에이스의 회사 소개에서 본 ``Make it Done!'' 문화도 이런 점에서 인상 깊었습니다. 저는 그 문장을 단순히 속도를 강조하는 표현으로 보기보다, 어려운 문제를 끝까지 기준 있게 정리하고 실제 결과로 연결하는 태도로 이해했습니다. 저 역시 프로젝트를 하면서 그와 비슷한 태도가 중요하다는 점을 배웠습니다.

\section*{5. 드림에이스에서 하고 싶은 일}
드림에이스에 지원하는 이유는 IVI, SDV, SILS를 서로 분리된 기술이 아니라 하나의 차량 소프트웨어 체계 안에서 다루고 있기 때문입니다. 특히 공식 사이트에서 확인한 IVI의 One-chip, Multi-OS 방향, SDV 영역의 virtualization과 AUTOSAR 기반 개발, SILS의 가상 ECU 기반 검증 플랫폼은 제가 앞으로 더 깊게 다루고 싶은 영역과 직접 맞닿아 있습니다.

저는 이번 프로젝트를 통해 차량 소프트웨어에서 중요한 것은 기능을 많이 붙이는 것이 아니라, 어떤 구조 위에 구현했고 어떤 기준으로 검증했는지를 끝까지 남기는 일이라는 점을 배웠습니다. 드림에이스에서도 저는 단순히 개발이나 시험 중 한 부분에 머무르지 않고, 구조를 이해한 상태에서 구현과 검증을 함께 잇는 엔지니어로 성장하고 싶습니다.

입사 후에는 드림에이스의 SILS·SDV 개발 흐름을 우선 빠르게 익히고, 프로젝트에서 쌓은 CANoe SIL, CAPL, 추적성, 시험 자산 정리 경험을 바탕으로, ``무엇이 동작하는가''뿐 아니라 ``왜 이 결과를 신뢰할 수 있는가''까지 설명할 수 있는 엔지니어가 되겠습니다. 이후에는 IVI 영역까지 이해를 넓혀, 구조와 검증을 함께 볼 수 있는 엔지니어로 성장하고 싶습니다.

\end{document}
